\documentclass[conference]{IEEEtran}
\IEEEoverridecommandlockouts
\usepackage{cite}
\usepackage{amsmath,amssymb,amsfonts}
\usepackage{algorithm}
\usepackage{algorithmic}
\usepackage{graphicx}
\usepackage{textcomp}
\usepackage{xcolor}
\usepackage{hyperref}
\usepackage{booktabs}
\usepackage{multirow}
\usepackage{threeparttable}
\def\BibTeX{{\rm B\kern-.05em{\sc i\kern-.025em b}\kern-.08em
    T\kern-.1667em\lower.7ex\hbox{E}\kern-.125emX}}

% ============================================================
% NOTE TO AUTHOR (delete before submission)
% Placeholders marked \NUM{...} must be filled from the re-run.
% ============================================================
\newcommand{\NUM}[1]{\textcolor{blue}{\texttt{[#1]}}}

\begin{document}

\title{Causal Influence Graph with Adaptive Mean Field: Uncertainty-Aware Core-Peripheral Modelling for Non-Stationary Multi-Agent Reinforcement Learning\\
\thanks{Identify applicable funding agency here. If none, delete this.}
}

\author{\IEEEauthorblockN{1\textsuperscript{st} Given Name Surname}
\IEEEauthorblockA{\textit{dept. name of organization (of Aff.)} \\
\textit{name of organization (of Aff.)}\\
City, Country \\
email address or ORCID}
\and
\IEEEauthorblockN{2\textsuperscript{nd} Given Name Surname}
\IEEEauthorblockA{\textit{dept. name of organization (of Aff.)} \\
\textit{name of organization (of Aff.)}\\
City, Country \\
email address or ORCID}
}

\maketitle

% ============================================================
\begin{abstract}
Multi-agent reinforcement learning (MARL) is destabilised by non-stationarity: because agents learn concurrently, the environment each agent faces keeps changing. Existing remedies compress the population by mean-field aggregation, weight neighbours by attention, or model opponents explicitly, but they conflate two distinct tiers of the problem---\emph{structural} non-stationarity (which agents materially influence a given ego-agent) and \emph{behavioural} non-stationarity (how those agents currently act).

We propose Causal Influence Graph with Adaptive Mean Field (CIG-AMF), which maintains an ego-centric directed influence graph and separates the two tiers by construction. Structural influence is defined as an interventional estimand and made identifiable on a randomised subset of transitions by injecting sparse forced actions ($\varepsilon$-forcing), which severs confounding mechanically and guarantees overlap; the budget is steered toward pairs the model is least certain about, subject to a floor that preserves overlap for every agent. The estimand is recovered by a doubly-robust estimator whose propensity is known exactly, since the behaviour policy is the learner's own; this yields \emph{signed}, multi-horizon influence estimates with ensemble uncertainty. Rather than collapsing each neighbour into one scalar, CIG-AMF assembles a six-dimensional \emph{influence signature} (sign, magnitude, uncertainty, temporal volatility, context-conditionality, latency) and uses it to route neighbours into memory slots with fixed functional meaning---\emph{beneficial}, \emph{harmful}, \emph{neutral}, \emph{anomalous}---plus free slots for residual structure. Core membership is decided by a lower-confidence-bound rule on bias-corrected beliefs updated on a Robbins--Monro schedule, which admits an almost-sure convergence argument; because that schedule also makes beliefs progressively rigid, a detected structural shift inflates uncertainty and reopens the learning rate, so the core contracts while the new regime is unfamiliar and re-expands as evidence accumulates.

We evaluate on an adaptive resource-flow arena under behavioural drift and structural shift, calibrate the estimator against clone-state interventions in a tiny-oracle companion, and report a \emph{selectivity ratio} that directly measures whether structural beliefs move under structural shift while staying still under behavioural drift. \NUM{fill: selectivity ratio, calibration Spearman, slot-usage entropy, reward and throughput deltas from the re-run}.
\end{abstract}

\begin{IEEEkeywords}
multi-agent reinforcement learning, non-stationarity, causal influence, doubly-robust estimation, mean field approximation, role discovery, structural belief, capacity allocation.
\end{IEEEkeywords}

% ============================================================
\section{Introduction}

Multi-agent reinforcement learning (MARL) underpins sequential decision-making in traffic control, resource allocation, distributed coordination, and strategic game-playing. Its defining difficulty is non-stationarity: because every agent learns concurrently, the transition and reward laws experienced by ego-agent $i$ depend on the evolving policies $\pi_{-i}$ of everyone else, violating the stationarity assumptions behind most RL algorithms \cite{yang2018mfmarl, busoniu2008survey, papoudakis2019nonstationarity}. We argue that this single label hides two tiers that demand different treatment. \emph{Structural} non-stationarity concerns \emph{which} neighbours materially affect $i$'s outcome---whether one controls a bottleneck, blocks a lane, or relays a signal---and is tied to task and environment structure. \emph{Behavioural} non-stationarity concerns \emph{how} those neighbours currently act as their policies drift. Conflating the tiers produces two failure modes: a structure assumed fixed cannot react to a genuinely moved bottleneck, whereas a structure re-estimated at behavioural speed mistakes transient policy noise for structural change, yielding oscillating neighbour sets and noisy policy inputs.

Existing lines each address part of this. Mean-field MARL makes many-agent learning tractable by replacing pairwise interaction with a population statistic \cite{yang2018mfmarl}, and weighted variants such as GAT-MF and adaptive mean field restore some heterogeneity through attention \cite{hao2023gatmf, wang2024amf}; but attention weights are estimated from observational trajectories and cannot separate structural importance from transient behavioural correlation. Causal mean-field Q-learning introduces intervention-inspired weights \cite{ma2023cmfq}, and a parallel line uses causal influence as an \emph{intrinsic reward} to promote coordination \cite{jaques2019influence, du2024scic, yu2026magic}; both improve the aggregate or the exploration signal, yet neither changes a neighbour's \emph{representational status}. Pieroth et al.\ formalise influence measurement functions with convergence guarantees, but deliberately avoid counterfactual actions, produce an unsigned magnitude, and state that using such measures to \emph{improve learning} remains future work \cite{pieroth2024influence}. Complementarily, coordination graphs sparsify interaction by payoff variance \cite{boehmer2020dcg, wang2022casec}, role-based MARL discovers agent roles from trajectories or action effects \cite{wang2020roma, wang2021rode, zeng2023sird, hu2024acorm}, and belief or meta-learning methods track latent task structure \cite{nguyen2011baipomdp, zintgraf2020varibad, zhang2024corep}. Roles in this literature are inferred from \emph{observed behaviour}, assigned \emph{globally}, and used to condition an agent's \emph{own} policy---so two vehicles that both sit still, one blocking the ego's lane and one parked harmlessly, receive the same role.

The gap is therefore specific. No existing method simultaneously (i) states structural influence as an interventional estimand and makes it identifiable during online training, (ii) estimates it with a \emph{signed}, uncertainty-aware, debiased estimator, (iii) represents each ego-neighbour relation as a multi-dimensional profile rather than one scalar, and (iv) uses that profile to allocate modelling capacity and to give memory slots fixed functional meaning. CIG-AMF fills exactly this gap, and its four contributions follow that order. First, we make the interventional target identifiable by \emph{$\varepsilon$-forcing}: with small probability a neighbour's action is drawn uniformly at random, severing confounding mechanically and guaranteeing overlap---a randomised experiment that costs nothing because the simulator is ours. Second, we estimate the target with a doubly-robust estimator whose propensity is \emph{exact}, since the behaviour policy is the learner's own; this is a structural advantage of MARL over observational causal inference, and it yields signed, multi-horizon estimates. Third, we replace the scalar weight with a six-dimensional influence signature and use it to route peripheral neighbours into slots with fixed semantics, regularised against the slot collapse that makes unconstrained soft routing degenerate to a mean. Fourth, we select core neighbours by a lower-confidence-bound rule on bias-corrected beliefs updated on a Robbins--Monro schedule, which both removes a saturation pathology and admits an almost-sure convergence argument in the sense of \cite{borkar2008sa, pieroth2024influence}.

Empirically, we report three findings that the design predicts. Calibrated against clone-state interventions, the signed doubly-robust estimator attains rank correlation \NUM{r\_signed} with true interventional effects versus \NUM{r\_range} for the unsigned magnitude used by prior influence measures, and sign agreement \NUM{s\_agree}. A \emph{selectivity ratio}---belief movement under structural shift divided by belief movement under behavioural drift---reaches \NUM{sel\_ratio} for CIG-AMF against \NUM{sel\_base} for correlation-based weighting, giving the first direct measurement that the two tiers are separated rather than merely postulated. Finally, semantically routed slots keep usage entropy at \NUM{ent} of maximum where unconstrained routing collapses to \NUM{ent0}, converting a previously null ablation into a measurable gain of \NUM{gain} at \NUM{cost} of the throughput of a single aggregate.

% ============================================================
\section{Background}

We isolate the two ideas needed to read the method: why aggregation discards the structure we want, and why observational influence weights are not the interventional quantity they are meant to approximate.

\subsection{Mean-Field Aggregation Discards Heterogeneous Influence}

Pairwise interaction modelling scales quadratically in the number of agents, and the joint action space grows exponentially. Mean-field MARL avoids both by factorising the pairwise value function and replacing individual neighbours with an aggregate statistic \cite{yang2018mfmarl}. Writing the mean neighbour action $\bar{a}^{i}=|\mathcal{N}_i|^{-1}\sum_{j\in\mathcal{N}_i}a^{j}$ and expanding to second order, the fluctuation terms $\delta a^{i,j}=a^{j}-\bar{a}^{i}$ sum to zero by construction, so the first-order term vanishes and
\begin{equation}
Q^{i}(s,\mathbf{a})\;\approx\;Q^{i}(s,a^{i},\bar{a}^{i}),
\label{eq:mf}
\end{equation}
reducing a population interaction to a two-body one. The approximation is exact only if neighbours are exchangeable contributors. That premise fails precisely in the cases we care about: a vehicle occupying a narrow lane matters more than many vehicles moving freely, and a station controller can dominate the reward trajectory of several agents at once. Attention-weighted variants \cite{hao2023gatmf, wang2024amf} restore heterogeneity by learning $\bar{a}^{i}=\sum_j w_{ij}a^{j}$, and coordination graphs sparsify interaction by payoff variance \cite{boehmer2020dcg, wang2022casec}. Both improve fidelity, but every neighbour remains inside the aggregation regime: the weights change, the representational status does not. CIG-AMF instead uses estimated influence as a \emph{partition} criterion, so that a small believed-influential set leaves the aggregate entirely and receives explicit pair-specific modelling.

\subsection{From Correlational Weights to an Interventional Estimand}

Weights learned from observed trajectories measure association. In MARL the two diverge for a concrete reason: a neighbour's action is a function of the shared state, so any statistic conditioned on $a_j$ absorbs the influence of common causes. Two vehicles braking at a distant red light co-vary strongly and influence each other not at all. The quantity we actually want is interventional---what would happen to $i$'s return \emph{if} $j$ acted differently, with the context held fixed.

The counterfactual difference-of-two-scenarios idea has a long lineage in MARL, but aimed at a different question. Difference rewards and the aristocrat utility marginalise an agent's \emph{own} action against a default or policy-averaged baseline to solve credit assignment \cite{wolpert2002difference, tumer2007aristocrat}, and COMA computes this baseline in a single forward pass of a centralised critic \cite{foerster2018coma}. Our estimand shares that algebraic form but reverses the target: we marginalise the \emph{neighbour's} action against $i$'s return, which asks who influences whom rather than who deserves credit. Recent causal MARL estimates interventional influence and converts it into an intrinsic reward, situation-dependently \cite{du2024scic} or over multiple steps with an advantage gate that recognises strong influence need not be useful influence \cite{yu2026magic}; the estimate steers \emph{what to do}, not \emph{whom to model}. Agent-specific effects give conditions under which agent-to-agent counterfactual effects are identifiable, but require structural assumptions such as noise monotonicity and are developed offline for accountability \cite{triantafyllou2024ase}.

Two consequences shape our method. First, because the difficulty is confounding rather than the formalism, we avoid strong structural assumptions and instead randomise: sparse forced actions restore the interventional distribution by construction, and known propensities enable doubly-robust estimation \cite{jiang2016dr, thomas2016hcope}. Second, because sign carries decision-relevant information that magnitude destroys---a neighbour who blocks and one who assists are opposites, not equals \cite{yu2026magic}---we retain it, whereas prior influence measurement is explicitly unsigned \cite{pieroth2024influence}.

% ============================================================
\section{Research Aim, Questions, and Hypotheses}

\subsection{Aim}
To develop and evaluate a MARL framework that (i) makes structural influence an identifiable interventional quantity during online training, (ii) uses signed, uncertainty-aware influence profiles to allocate modelling capacity between explicitly modelled and compressed neighbours, and (iii) demonstrably separates structural from behavioural non-stationarity.

\subsection{Research Questions}
\begin{itemize}
\item \textbf{RQ1 (Identification).} Do sparse randomised interventions and doubly-robust estimation produce influence estimates that agree with controlled clone-state interventions in sign and rank, and does the \emph{signed} estimand outperform the unsigned magnitude used by prior influence measures?
\item \textbf{RQ2 (Tier separation).} Do structural beliefs remain stable under behavioural drift while responding under genuine structural shift, and does this selectivity exceed that of correlation-based weighting?
\item \textbf{RQ3 (Capacity allocation).} Does routing peripheral neighbours by influence signature into semantically fixed slots preserve more decision-relevant information than a single aggregate, at a cost below explicit modelling---and does it avoid the slot collapse that makes unconstrained routing degenerate?
\end{itemize}

\subsection{Hypotheses}
\begin{itemize}
\item \textbf{H1.} Under $\varepsilon$-forcing, the doubly-robust signed estimator will show lower bias and higher rank correlation against oracle interventions than a plug-in estimator, and higher correlation than an unsigned range statistic; bias will decrease with $\varepsilon$.
\item \textbf{H2.} The selectivity ratio of CIG-AMF will exceed unity by a clear margin, whereas correlation-based weighting will be close to unity, because association shifts under both conditions.
\item \textbf{H3.} Semantic slot routing with load-balancing and orthogonality regularisation will retain usage entropy near maximum and improve value-prediction agreement with an explicit reference relative to a single aggregate, while unconstrained soft routing will collapse and show no such gain.
\end{itemize}

% ============================================================
\section{Methodology}

\subsection{Setting and Notation}

We consider $N$ agents in a partially observable stochastic game. At step $t$, agent $i$ receives observation $o_i^{(t)}$, takes action $a_i^{(t)}\in\mathcal{A}$, and receives reward $r_i^{(t)}$. Its $h$-step discounted return is
\begin{equation}
R_i^{(h)}(t)=\sum_{k=0}^{h-1}\gamma^{k}\,r_i(t+k), \qquad h=1,\dots,H .
\label{eq:return}
\end{equation}
For each ego-agent $i$ we maintain an ego-centric directed graph over $\mathcal{N}_i$, partitioned online into an explicitly modelled \emph{core} $\mathcal{C}_i$ and a compressed \emph{periphery} $\mathcal{P}_i$, with $\mathcal{C}_i\cup\mathcal{P}_i=\mathcal{N}_i$ and $\mathcal{C}_i\cap\mathcal{P}_i=\emptyset$. Core neighbours are represented by pair-specific latents $z_{ij}$ pooled into $Z_i$; peripheral neighbours are compressed into $M_i$; belief statistics are summarised in $B_i$. Policy and value share the conditioning
\begin{equation}
\pi_i(\cdot\mid o_i,Z_i,M_i,B_i),\qquad V_i(o_i,Z_i,M_i,B_i).
\label{eq:policy}
\end{equation}

\subsection{Structural Influence: Estimand and Identification}

\subsubsection{Estimand}
Following the aristocrat-utility form \cite{tumer2007aristocrat} but reversing the target from own action to neighbour action, we define the structural influence of $j$ on $i$ in state $s$ as the interventional contrast between $j$'s realised action and $j$'s own policy-averaged behaviour:
\begin{equation}
\begin{aligned}
w_{ij}(s,a_j) \;=\; &\mathbb{E}\!\left[R_i^{(H)}\mid \mathrm{do}(a_j),\,s\right] \\
&-\; \mathbb{E}_{a'\sim\pi_j(\cdot\mid s)}\!\left[\mathbb{E}\!\left[R_i^{(H)}\mid \mathrm{do}(a'),\,s\right]\right].
\end{aligned}
\label{eq:estimand}
\end{equation}
The sign is meaningful and is retained throughout: $w_{ij}>0$ means $j$'s realised action helped $i$ relative to $j$'s typical behaviour; $w_{ij}<0$ means it harmed $i$. The policy-averaged baseline, rather than an arbitrarily chosen alternative action, follows \cite{foerster2018coma} and reduces variance because it averages over the full action distribution. We also define the unsigned range statistic $u_{ij}(s)=\max_a \mathbb{E}[R_i^{(H)}\mid\mathrm{do}(a),s]-\min_a \mathbb{E}[R_i^{(H)}\mid\mathrm{do}(a),s]$, which corresponds to the impact sample of \cite{pieroth2024influence} and which we retain as a baseline.

\subsubsection{Assumptions}
\begin{itemize}
\item[\textbf{A1}] \emph{Bounded local influence and horizon.} $j$'s action affects $i$ directly or through short local chains, with practically relevant effect within horizon $H$. This matches the arena, where lanes, stations, and demand zones define bounded interaction regions.
\item[\textbf{A2}] \emph{Overlap (positivity).} Every action retains positive probability under the behaviour policy: $b_j(a\mid s)\ge \varepsilon/|\mathcal{A}|>0$ for all $a$. Without A2 the counterfactual action $a'$ may lie outside the data support and the outcome model must extrapolate.
\item[\textbf{A3}] \emph{Randomised intervention.} At each step, independently with probability $\varepsilon_j\ge\varepsilon_{\min}>0$, agent $j$'s action is drawn uniformly from $\mathcal{A}$ independently of $s$, $a_i$, and $a_{-j}$. The \emph{selection} probability $\varepsilon_j$ may depend on the learner's own state, but the forced action itself must not.
\item[\textbf{A4}] \emph{Timescale separation.} Structural dependencies change more slowly than behaviour, motivating the two-timescale schedule.
\end{itemize}

\subsubsection{$\varepsilon$-forcing and why it matters}
A2 and A3 are not assumptions we hope hold; we \emph{enforce} them. Because the environment is a simulator we control, the learner may occasionally override a neighbour's sampled action. Under A3 the forced action is independent of every common cause, so on forced steps
\begin{equation}
\mathbb{E}\!\left[R_i^{(H)}\mid \mathrm{do}(a_j=a),\,s\right]
=\mathbb{E}\!\left[R_i^{(H)}\mid a_j=a,\,s,\,F_j=1\right],
\label{eq:identification}
\end{equation}
where $F_j$ indicates forcing. That is, on the forced sub-population the interventional quantity of \eqref{eq:estimand} reduces to an ordinary conditional expectation, identified without any structural assumption on the noise---in contrast to identification routes requiring noise monotonicity or bijective mechanisms \cite{triantafyllou2024ase}. Two qualifications are essential and we state them rather than leave them implicit. First, the guarantee is \emph{local to the randomised subset}: with $\varepsilon=0.03$ only three percent of transitions carry it, while the outcome model is trained on the full distribution and therefore remains exposed to residual confounding, which the correction of \eqref{eq:dr} mitigates but does not eliminate. Second, identification is purchased with perturbation, and we report its cost in return rather than assume it negligible.

The effective behaviour policy is the mixture
\begin{equation}
b_j(a\mid s)=\varepsilon_j\,\tfrac{1}{|\mathcal{A}|}+(1-\varepsilon_j)\,\pi_j(a\mid s),
\label{eq:behaviour}
\end{equation}
which is \emph{known exactly}, since $\pi_j$ is the learner's own network and $\varepsilon_j$ is a chosen quantity. Off-policy evaluation in observational domains must estimate this and inherits its error; in MARL it is available in closed form, an advantage we exploit next.

\subsubsection{Allocating the intervention budget}
A uniform $\varepsilon$ spends most of the budget on pairs already well understood. Since the budget is paid in return, we allocate it where it buys the most information, setting
\begin{equation}
\varepsilon_j \;\propto\; \varepsilon_{\min} + (1-\varepsilon_{\min})\,\frac{\hat\sigma_{\cdot j}}{\overline{\hat\sigma}},
\qquad
\tfrac{1}{N}\textstyle\sum_j \varepsilon_j=\varepsilon,
\label{eq:targeted}
\end{equation}
where $\hat\sigma_{\cdot j}$ aggregates epistemic uncertainty over egos observing $j$, and the normalisation preserves the total budget. Two constraints make this admissible. The floor $\varepsilon_{\min}>0$ preserves A2 for every agent: an agent never forced would lose overlap entirely, and every causal statement about it would lose its basis. And because the forced action remains uniform and independent of the state given that forcing occurred, \eqref{eq:identification} is unaffected---the \emph{selection} of whom to perturb may depend on anything, provided the realised $\varepsilon_j$ is recorded and enters \eqref{eq:behaviour}, which it does.

\subsection{Doubly-Robust Signed Proxy Estimation}

\subsubsection{Outcome model}
An ensemble of $K$ networks $f_{\theta_k}$ predicts the ego return for a hypothetical neighbour action, conditioned on the ego state and action, the target neighbour's action, and the remaining core and peripheral context with $j$ removed:
\begin{equation}
f_{\theta_k}\!\left(o_i,a_i,a,\,Z_i^{-j},\,M_i^{-j}\right)\;\longmapsto\;
\left(\hat{R}_i^{(1)},\dots,\hat{R}_i^{(H)}\right).
\label{eq:outcome}
\end{equation}
Two design choices differ from a naive conditioning set. First, the belief summary $B_i$ is \emph{excluded} from \eqref{eq:outcome}: $B_i$ is computed from the very influence estimates this model produces, so including it creates a feedback loop in which the model can confirm its own partition---a confounder introduced by the architecture rather than the environment. Second, the head is multi-horizon: a blocker acts immediately whereas a relay's benefit appears several steps later, and a single aggregated return cannot distinguish them.

\subsubsection{Doubly-robust contrast}
Let $R_i^{(H)}$ be the realised return and $a_j$ the realised neighbour action. Define the augmented outcome
\begin{equation}
\psi_{ij}(a)=f_{\theta}(a)+\frac{\mathbb{I}[a_j=a]}{b_j(a\mid s)}\Big(R_i^{(H)}-f_{\theta}(a_j)\Big),
\label{eq:psi}
\end{equation}
where $f_\theta(a)$ abbreviates \eqref{eq:outcome} evaluated at neighbour action $a$. Substituting \eqref{eq:psi} into \eqref{eq:estimand} and using $\sum_a \pi_j(a)\,\mathbb{I}[a_j=a]/b_j(a)=\pi_j(a_j)/b_j(a_j)$ gives the closed form
\begin{equation}
\begin{aligned}
\hat{w}_{ij}=\;&\underbrace{f_{\theta}(a_j)-\textstyle\sum_{a}\pi_j(a\mid s)\,f_{\theta}(a)}_{\text{plug-in contrast}}\\
&+\underbrace{\Big(R_i^{(H)}-f_{\theta}(a_j)\Big)\left(\frac{1}{b_j(a_j\mid s)}-1\right)}_{\text{doubly-robust correction}} .
\end{aligned}
\label{eq:dr}
\end{equation}
The two terms explain the estimator's robustness. If the outcome model is correct, the residual $R_i^{(H)}-f_\theta(a_j)$ has zero conditional mean and the correction vanishes, leaving the exact plug-in contrast. If the outcome model is biased, the correction re-weights the residual by the inverse propensity and removes the first-order error. The estimator is therefore unbiased when \emph{either} component is correct \cite{jiang2016dr}; because \eqref{eq:behaviour} is exact, one component is correct by construction. The importance ratio is clipped at $c_{\mathrm{iw}}$ to bound variance \cite{thomas2016hcope}. For comparison we retain the plug-in variant (correction disabled) and the unsigned range statistic as ablations.

\subsubsection{Ensemble uncertainty}
Each member $k$ is trained on an independent bootstrap subsample with its own initialisation, and interventional ($F_j=1$) samples are oversampled by a factor $\eta$ because they carry identification content. Writing $\hat{w}^{(k)}_{ij}$ for member $k$'s estimate,
\begin{equation}
\mu_{ij}=\frac{1}{K}\sum_{k}\hat{w}^{(k)}_{ij},
\qquad
\sigma_{ij}=\sqrt{\frac{1}{K-1}\sum_{k}\big(\hat{w}^{(k)}_{ij}-\mu_{ij}\big)^{2}} .
\label{eq:ensemble}
\end{equation}
Independent subsampling is essential: members trained on identical batches converge to nearly identical functions and $\sigma_{ij}$ degenerates to a deterministic schedule rather than an epistemic quantity \cite{lakshminarayanan2017ensembles}. We report ensemble disagreement as a diagnostic so that this failure is visible rather than silent.

\subsection{Structural Belief and Core Selection}

\subsubsection{Bias-corrected Robbins--Monro smoothing}
For each directed pair we maintain $b_{ij}=(\bar\mu_{ij},\bar\sigma_{ij})$ by exponential smoothing,
\begin{equation}
\bar\mu^{(t)}_{ij}=(1-\alpha^{(t)}_{ij})\bar\mu^{(t-1)}_{ij}+\alpha^{(t)}_{ij}\mu^{(t)}_{ij},
\label{eq:smooth_mu}
\end{equation}
\begin{equation}
\bar\sigma^{(t)}_{ij}=(1-\alpha^{(t)}_{ij})\bar\sigma^{(t-1)}_{ij}+\alpha^{(t)}_{ij}\sigma^{(t)}_{ij},
\label{eq:smooth_sigma}
\end{equation}
with the step size
\begin{equation}
\alpha^{(t)}_{ij}=\frac{\alpha_0}{t^{d}\,\bigl(1+c_\sigma\,\sigma^{(t)}_{ij}\bigr)},
\qquad d\in(\tfrac12,1].
\label{eq:alpha}
\end{equation}
The factor $(1+c_\sigma\sigma)^{-1}$ makes high-uncertainty evidence less disruptive. The factor $t^{-d}$ is what makes the schedule admissible: it gives $\sum_t\alpha_t=\infty$ and $\sum_t\alpha_t^2<\infty$, the Robbins--Monro conditions required by stochastic-approximation arguments \cite{borkar2008sa}, and used in the same iteration form by \cite{pieroth2024influence}. A step size bounded away from zero, as in fixed-rate smoothing, violates the square-summability condition and admits no such statement.

\smallskip\noindent\textbf{Proposition 1 (convergence).} \emph{Under A4, if the proxy sequence $\mu^{(t)}_{ij}$ converges almost surely and $\alpha^{(t)}_{ij}$ satisfies \eqref{eq:alpha} with $d\in(\tfrac12,1]$, then $\bar\mu^{(t)}_{ij}$ converges almost surely to the limit of $\mu^{(t)}_{ij}$.} The argument is the standard stochastic-approximation one \cite{borkar2008sa}: \eqref{eq:smooth_mu} is linear in its first argument, the driving process is ergodic, and the associated ODE $\dot{\bar\mu}=\mu^{\ast}-\bar\mu$ has a unique globally asymptotically stable equilibrium.

Decaying step sizes have a side effect that must be corrected. After $t$ updates the accumulated weight on evidence is $1-\Pi^{(t)}_{ij}$ with $\Pi^{(t)}_{ij}=\prod_{s\le t}(1-\alpha^{(s)}_{ij})$, so the remainder still belongs to the initial values $\mu_0=0$ and $\sigma_0=1$. Left uncorrected, $\bar\mu$ is shrunk toward zero and---more damagingly---$\bar\sigma$ retains part of the maximal-uncertainty prior and is inflated. We therefore use the debiased quantities
\begin{equation}
\hat\mu_{ij}=\frac{\bar\mu_{ij}-\Pi_{ij}\mu_0}{1-\Pi_{ij}},
\qquad
\hat\sigma_{ij}=\frac{\bar\sigma_{ij}-\Pi_{ij}\sigma_0}{1-\Pi_{ij}},
\label{eq:debias}
\end{equation}
which generalise Adam's bias correction \cite{kingma2015adam} to a non-zero initialisation and to a step size that varies with observed uncertainty. All decisions below use \eqref{eq:debias}; the raw values are logged only for diagnostics.

\subsubsection{Controlled forgetting on detected shift}
The schedule of \eqref{eq:alpha} creates a tension the paper must resolve rather than ignore. Square-summability is what licenses Proposition 1, but it also makes beliefs progressively rigid: after many updates $\alpha$ is small, and a genuinely changed dependency needs a long time to overturn an entrenched estimate. In a stationary problem that rigidity is the desired behaviour; in the non-stationary problem this paper addresses, it is a defect.

When a shift is detected (\S\ref{sec:trigger}), we therefore re-anchor and inflate. Writing $\Pi_{ij}$ for the accumulated product in \eqref{eq:debias}, the affected pairs are reset to
\begin{equation}
\mu_0\!\leftarrow\!\hat\mu_{ij},
\quad
\sigma_0\!\leftarrow\!\min(\nu\,\hat\sigma_{ij},\,\sigma_{\max}),
\quad
\Pi_{ij}\!\leftarrow\!1,
\quad
t\!\leftarrow\!1,
\label{eq:inflate}
\end{equation}
with inflation factor $\nu>1$. The re-anchoring is what distinguishes this from a reset: the current debiased estimate becomes the new prior, so evidence is not discarded, only its \emph{credibility} is reduced. Setting $t\!\leftarrow\!1$ reopens the learning rate so that new evidence can act.

The mechanism has a consequence we did not design but which follows from \eqref{eq:lcb}: inflating $\hat\sigma$ lowers every $g_{ij}$, so fewer neighbours clear $\tau_{\mathrm{in}}$ and \emph{the core contracts automatically}. The system becomes selective precisely while the new regime is unfamiliar, then re-expands as evidence accumulates---without any hand-written rule. This is covariance inflation in the Kalman sense, and is of a piece with discounting and sliding windows in non-stationary bandits.

\subsubsection{Lower-confidence-bound core selection}
Core membership is decided by the uncertainty-penalised score
\begin{equation}
g_{ij}=\big|\hat\mu_{ij}\big|-\kappa\,\hat\sigma_{ij},
\label{eq:lcb}
\end{equation}
with hysteresis to prevent oscillation near the boundary:
\begin{equation}
j\in\mathcal{C}^{(t)}_i \iff
\begin{cases}
g_{ij}>\tau_{\mathrm{in}} & \text{if } j\notin\mathcal{C}^{(t-1)}_i,\\[2pt]
g_{ij}>\tau_{\mathrm{out}} & \text{if } j\in\mathcal{C}^{(t-1)}_i,
\end{cases}
\label{eq:hysteresis}
\end{equation}
where $\tau_{\mathrm{out}}<\tau_{\mathrm{in}}$. Placing $\hat\sigma$ additively is deliberate. A score of the form $\mathrm{sigmoid}\big((|\hat\mu|-\tau)/\hat\sigma\big)$ divides by the uncertainty and therefore saturates: when $\hat\sigma$ is small the argument is amplified and the score becomes a hard $0/1$ regardless of the influence magnitude, whereas when $\hat\sigma$ is large all neighbours compress toward a common value. In both regimes the \emph{scale} of $\hat\sigma$, not the structural signal, determines the partition. The additive penalty of \eqref{eq:lcb} keeps the ordering governed by $\hat\mu$ while still discounting uncertain pairs, and is the standard confidence-bound construction used for selection under uncertainty; it is also consistent with variance-based edge sparsification in coordination graphs \cite{wang2022casec}.

\subsubsection{Adaptive core budget}
Rather than a fixed cap, the target core size adapts to how concentrated influence is. With $p_{ij}=|\hat\mu_{ij}|/\sum_{l}|\hat\mu_{il}|$ and normalised entropy $\tilde{h}_i=-\sum_j p_{ij}\log p_{ij}/\log|\mathcal{N}_i|$,
\begin{equation}
k_i=\big\lceil k_{\min}+\tilde{h}_i\,(k_{\max}-k_{\min}) \big\rceil .
\label{eq:budget}
\end{equation}
Concentrated influence yields a small core; diffuse influence yields a larger one. We report the fraction of updates at which $k_i$ binds at $k_{\max}$, because a capacity allocation that always saturates is a constant, not an allocation.

\subsection{Influence Signatures and Role-Structured Peripheral Memory}

\subsubsection{Signature}
Compressing a relation into one scalar discards structure that decisions depend on. A blocker exerts a strongly negative, immediate effect confined to specific contexts; a relay exerts a positive effect that appears several steps later across many contexts; a resource competitor exerts a weak negative effect uniformly. These can share an average magnitude while requiring different responses. We therefore summarise each directed pair by
\begin{equation}
\phi_{ij}=\big(\hat\mu_{ij},\;\overline{|w|}_{ij},\;\hat\sigma_{ij},\;v^{\mathrm{tim}}_{ij},\;v^{\mathrm{ctx}}_{ij},\;\ell_{ij}\big)\in\mathbb{R}^{6},
\label{eq:signature}
\end{equation}
whose components are the signed influence; the mean absolute influence $\overline{|w|}_{ij}$; the epistemic uncertainty; the temporal standard deviation of $w$ over a window; the standard deviation of context-conditional means across context keys; and the latency
\begin{equation}
\ell_{ij}=\frac{\sum_{h=1}^{H}(h-1)\,\big|\hat{w}^{(h)}_{ij}\big|}{\sum_{h=1}^{H}\big|\hat{w}^{(h)}_{ij}\big|+\epsilon},
\label{eq:latency}
\end{equation}
the centre of mass of the effect over the horizon. The first two components are deliberately distinct: $|\hat\mu|$ measures net directed influence whereas $\overline{|w|}$ measures magnitude irrespective of direction, and their ratio separates a consistently harmful neighbour from one whose influence reverses. Components four and five separate temporal instability from genuine context-conditionality, the latter being the situation-dependence emphasised by \cite{du2024scic, seitzer2021cai}.

\subsubsection{Semantic slot routing}
Unconstrained soft routing---a softmax head trained only through the distant policy loss---has no signal telling any slot what to hold, and degenerates in one of two ways: a single slot absorbs all mass, or all slots receive near-uniform mass and each reproduces the global mean. Either way the multi-slot memory becomes a copy of a single aggregate. This is the well-documented collapse of mixture-of-experts routing and slot attention \cite{fedus2022switch, locatello2020slotattention}. We therefore give slots fixed functional meaning derived from \eqref{eq:signature}. With gates
\begin{align}
g^{\mathrm{anom}}_{ij}&=\mathrm{sigmoid}\!\big(\tfrac{\varrho}{\sigma_{\mathrm{hi}}}(\hat\sigma_{ij}-\sigma_{\mathrm{hi}})\big), \label{eq:gate_anom}\\
g^{+}_{ij}&=\mathrm{sigmoid}\!\big(\tfrac{\varrho}{\tau_\rho}(\hat\mu_{ij}-\tau_\rho)\big),\quad
g^{-}_{ij}=\mathrm{sigmoid}\!\big(\tfrac{\varrho}{\tau_\rho}(-\hat\mu_{ij}-\tau_\rho)\big), \label{eq:gate_pm}
\end{align}
and $g^{0}_{ij}=\mathrm{clip}(1-g^{+}_{ij}-g^{-}_{ij},0,1)$, the semantic assignment is
\begin{equation}
\rho_{ij}\propto\Big[(1-g^{\mathrm{anom}}_{ij})g^{+}_{ij},\;(1-g^{\mathrm{anom}}_{ij})g^{-}_{ij},\;(1-g^{\mathrm{anom}}_{ij})g^{0}_{ij},\;g^{\mathrm{anom}}_{ij}\Big],
\label{eq:rho}
\end{equation}
normalised to sum to one, with slots ordered \emph{beneficial}, \emph{harmful}, \emph{neutral}, \emph{anomalous}. Two details matter. The gate sharpness $\varrho$ is divided by the threshold it acts on, making it dimensionless; an unnormalised sharpness leaves the transition width dependent on the reward scale and mis-assigns near-zero influence. The anomalous gate is applied first: when a relation is not yet understood, assigning it a valence is unwarranted, so uncertainty becomes a semantic category rather than only a control parameter. Thresholds $\tau_\rho,\sigma_{\mathrm{hi}}$ are set by percentiles of the observed distribution after warm-up, since the scale of $w$ is set by the environment's reward scale.

$K_f$ additional free slots capture residual structure. Their router is conditioned on $\rho_{ij}$,
\begin{equation}
u_{ij}=\mathrm{softmax}\!\big(W_f\,[\,h_{ij};\,\rho_{ij}\,]\big),
\qquad
\alpha_{ij}=\big[\tfrac12\rho_{ij};\;\tfrac12 u_{ij}\big],
\label{eq:free_slots}
\end{equation}
so that free capacity specialises \emph{within} a valence group rather than competing globally. Global partitioning is dominated by strongly influential relations, which compresses weak ones together; conditioning on valence first avoids this.

\subsubsection{Pooling and regularisation}
Peripheral item features are $x^{P}_{ij}=[a_j;\xi_{ij};\hat\mu_{ij};\hat\sigma_{ij}]$ with $\xi_{ij}$ the relative geometry, encoded to $h_{ij}$. The confidence weight and slot contents are
\begin{equation}
\beta_{ij}=\frac{(\beta_0+p_{ij})\big(|\hat\mu_{ij}|+\mu_0\big)}{1+\hat\sigma_{ij}},
\label{eq:beta}
\end{equation}
\begin{equation}
m_{iq}=\frac{\sum_{j\in\mathcal{P}_i}\beta_{ij}\,\alpha_{ij,q}\,h_{ij}}{\sum_{j\in\mathcal{P}_i}\beta_{ij}\,\alpha_{ij,q}+\epsilon},
\qquad
M_i=W_o\,\big[m_{i1};\dots;m_{iK}\big].
\label{eq:slots}
\end{equation}
The pooling in \eqref{eq:slots} is permutation-invariant in the sense of Deep Sets \cite{zaheer2017deepsets}; more expressive set encoders \cite{lee2019settransformer} could replace it, at higher cost. Two auxiliary losses guard against the two collapse modes. A load-balancing term, with $f_q$ the fraction of items whose arg-max slot is $q$ and $P_q$ the mean routing probability to $q$ \cite{fedus2022switch}, penalises monopolisation:
\begin{equation}
\mathcal{L}_{\mathrm{lb}}=K\sum_{q=1}^{K}f_q\,P_q .
\label{eq:lb}
\end{equation}
An orthogonality term penalises slots that hold interchangeable content, which load balancing alone does not detect:
\begin{equation}
\mathcal{L}_{\mathrm{orth}}=\frac{1}{K(K-1)}\sum_{q\neq r}\cos^{2}\!\big(m_{iq},\,m_{ir}\big),
\qquad
\mathcal{L}_{\mathrm{aux}}=\lambda_{\mathrm{lb}}\mathcal{L}_{\mathrm{lb}}+\lambda_{\mathrm{orth}}\mathcal{L}_{\mathrm{orth}} .
\label{eq:orth}
\end{equation}
This use of mutual-information-style specialisation pressure parallels role learning \cite{wang2020roma}, but the signal is interventional influence rather than observed trajectories, and it structures the representation \emph{of others} rather than the agent's own policy.

\subsection{Ego-Conditioned Pair-Relational Latents}

Core neighbours are represented by a recurrent pair latent
\begin{equation}
z^{(t)}_{ij}=f_z\big(z^{(t-1)}_{ij},\,o_i^{(t)},a_i^{(t)},\,o_j^{(t)},a_j^{(t)},\,\xi^{(t)}_{ij}\big),
\qquad Z_i=\mathrm{Pool}\{z_{ij}\}_{j\in\mathcal{C}_i},
\label{eq:zij}
\end{equation}
with a low-dimensional shadow state maintained for all neighbours and projected into $z_{ij}$ on promotion, avoiding cold starts. Training only on next-action prediction, $\mathcal{L}_{\mathrm{bc}}=-\log p(a^{(t+1)}_j\mid z^{(t)}_{ij})$, does not make the latent ego-specific: predicting $a_j$ requires no information about $i$, so gradients permit the encoder to ignore $o_i,a_i$ and converge to a global opponent model shared across egos. We therefore add two terms. An influence head regresses the signed multi-horizon influence, a quantity that depends on both agents:
\begin{equation}
\mathcal{L}_{w}=\mathrm{Huber}\Big(g_w(z_{ij}),\;\big[\hat{w}^{(1)}_{ij},\dots,\hat{w}^{(H)}_{ij}\big]\Big).
\label{eq:lw}
\end{equation}
A contrastive term separates the same neighbour viewed by different egos, with positives drawn from the same pair at other times and negatives from the same $j$ under different $i$:
\begin{equation}
\mathcal{L}_{\mathrm{con}}=-\log\frac{\sum_{p\in\mathcal{P}os}\exp(\langle \tilde{z},\tilde{z}_p\rangle/\varsigma)}{\sum_{c\in\mathcal{P}os\cup\mathcal{N}eg}\exp(\langle \tilde{z},\tilde{z}_c\rangle/\varsigma)},
\label{eq:lcon}
\end{equation}
with $\tilde{z}$ the $\ell_2$-normalised projection and $\varsigma$ a temperature. The negative set is restricted to the same neighbour across egos because that is precisely the invariance we wish to break: $j$ may block $i$ while assisting $i'$. The total is $\mathcal{L}_z=\mathcal{L}_{\mathrm{bc}}+\lambda_w\mathcal{L}_w+\lambda_c\mathcal{L}_{\mathrm{con}}$. We report a pair-specificity diagnostic---mean cosine similarity of $z_{ij}$ across egos for fixed $j$, relative to similarity across neighbours for fixed $i$---so that the claim is measured rather than asserted.

\subsection{Two-Timescale Schedule and Shift Detection}
\label{sec:trigger}

Policy, value, pair latents, shadow states, and the peripheral encoder update on the fast timescale. The proxy ensemble, structural beliefs, and the core-peripheral partition update on the slow timescale. Training proceeds in two stages: a seeded warm-up in which a weak prior (local distance, zone co-membership) provides a non-empty core while the replay buffer fills, and a learned-belief stage in which \eqref{eq:lcb}--\eqref{eq:hysteresis} take over.

\subsubsection{Why a learning detector cannot detect}
The natural trigger---the proxy's own prediction error---fails for two reasons that compound. A model still being trained quietly absorbs the very change it is meant to flag, so its error returns to baseline and the alarm never sounds. Worse, the proxy's inputs $Z^{-j}_i$ and $M^{-j}_i$ are functions of the current partition, so re-partitioning alone perturbs the error and re-triggers the mechanism: the detector responds to its own motion rather than to the environment.

We therefore monitor with a \emph{frozen, context-blind probe}. A small network is trained to predict $R_i^{(H)}$ from $(o_i,a_i,a_j)$ alone, then its weights are frozen:
\begin{equation}
q_\vartheta:\;(o_i,a_i,a_j)\longmapsto \hat{R}^{(1..H)}_i,
\qquad
\vartheta\;\text{fixed after snapshot},
\label{eq:probe}
\end{equation}
and scored on the most recent transitions. Freezing keeps its expectations anchored to the regime in which it was trained, so it cannot adapt away the signal. Withholding $Z,M,B$ makes it structurally incapable of responding to our own bookkeeping. Its residual rises when the environment's law changes, and for no other reason. Because the probe becomes permanently stale after a genuine shift, it is re-snapshotted once adaptation has settled; without this it would alarm indefinitely.

\subsubsection{A second, faster detector}
The probe inherits a delay: $R_i^{(H)}(t)$ is observable only at $t+H$. We therefore add an independent signal on the influence matrix $W$ itself, using the normalised movement of \eqref{eq:selectivity}. The justification is borrowed rather than assumed: influence measurement functions are continuous in the policy parameters \cite{pieroth2024influence}, so behavioural drift---a smooth motion of $\pi$---can only move $W$ smoothly. A discontinuity in $W$ therefore cannot be behavioural, and must be structural. This is the formal counterpart of the tier separation the paper argues for.

Both signals are standardised into $z$-scores over a trailing window---a change-detection convention dating to sequential inspection \cite{page1954cusum}---which makes the detection threshold dimensionless and transferable across reward scales, and either may fire (a conjunctive rule is available as an option). A refractory period follows each firing: without it, the partition change and belief inflation triggered by one detection would immediately re-trigger the next. On firing, the slow update rate is multiplied, uncertainty is inflated via \eqref{eq:inflate}, and the probe is scheduled for re-snapshot. Triggering on these quantities rather than on critic loss is deliberate, since critic error rises for many reasons unrelated to structure \cite{zhang2024corep}.

\subsection{Algorithm}

\begin{algorithm}[t]
\caption{CIG-AMF (one episode)}
\label{alg:cigamf}
\begin{algorithmic}[1]
\STATE Reset environment; obtain $\{o_i\}$
\FOR{each step $t$}
  \STATE For each $i$: build $Z_i$ from $\{z_{ij}\}_{j\in\mathcal{C}_i}$, $M_i$ via \eqref{eq:rho}--\eqref{eq:slots}, and $B_i$
  \STATE Sample $a_i\sim\pi_i(\cdot\mid o_i,Z_i,M_i,B_i)$; record $\pi_j(\cdot\mid s)$
  \STATE \textbf{$\varepsilon$-forcing:} with prob.\ $\varepsilon_j$ from \eqref{eq:targeted} override $a_j\sim\mathrm{Unif}(\mathcal{A})$; record $F_j$ and $b_j$ via \eqref{eq:behaviour}
  \STATE Step environment; cache $Z^{-j}_i,M^{-j}_i,B_i$ at action time
  \STATE Update $z_{ij}$ and shadow states \COMMENT{fast}
\ENDFOR
\STATE Compute $R_i^{(1..H)}$; push transitions with $b_j$, $F_j$, context key
\STATE Update $\pi,V$ with $\mathcal{L}_{\mathrm{aux}}$; update $\mathcal{L}_z$ \COMMENT{fast}
\IF{slow-timescale step}
  \STATE Train ensemble \eqref{eq:outcome} on per-member bootstrap batches
  \STATE Score pairs by \eqref{eq:dr}, \eqref{eq:ensemble}; update signatures \eqref{eq:signature}
  \STATE Update beliefs \eqref{eq:smooth_mu}--\eqref{eq:debias}; re-partition \eqref{eq:lcb}--\eqref{eq:budget}
  \STATE Warm-start $z_{ij}$ for promoted $j$
\ENDIF
\STATE Score frozen probe \eqref{eq:probe} and matrix movement \eqref{eq:selectivity}; if either $z$-score fires and not refractory: accelerate, inflate \eqref{eq:inflate}, schedule re-snapshot
\STATE After warm-up (once): calibrate $\tau_\rho,\sigma_{\mathrm{hi}}$ by percentile
\end{algorithmic}
\end{algorithm}

% ============================================================
\section{Experiments}

The evaluation is organised so that each experiment tests one claim, and so that a prerequisite is checked before the claims that depend on it.

\subsection{Environments}

The main benchmark is an adaptive resource-flow arena with multiple zones, each containing a processing chain of sources, buffers, lanes, stations, and sinks. Agents move, collect, transfer, process, deliver, toggle lanes, signal, or inspect local rules. Structured interdependence arises from bottlenecks that create high-influence agents, hidden rules that alter which agents matter, role heterogeneity, partial observability through local crops, congestion with delayed reward, and indirect effects through observable signals. The design follows benchmark lessons on social generalisation, stochasticity, and many-agent persistence \cite{leibo2021meltingpot, ellis2023smacv2, suarez2019neuralmmo, carroll2019overcooked}, while remaining controlled enough for clone-state intervention.

A tiny-oracle companion shares the ontology at smaller scale and supports state cloning, restoration, state-bank sampling, and short-horizon intervention rollout. It is used exclusively for calibration. Crucially, oracle and estimator must target the same quantity: the oracle returns a \emph{signed} contrast, so calibration uses the oracle-matched signed mode of the estimator; comparing a signed oracle against an unsigned estimator would be uninformative regardless of estimator quality.

\subsection{Experiment 0: Structure Sensitivity (prerequisite)}

Before asking whether a method recovers structure, we measure how much recovering structure is worth in this environment. We compare pure mean field, an \emph{oracle-core} variant given the ground-truth core at no learning cost, and full explicit local modelling. The gap between oracle-core and pure mean field upper-bounds the achievable return of any structural method here. If that gap is small relative to reward noise, no algorithmic change can be demonstrated, and the environment---not the method---must be revised by tightening lanes, increasing congestion penalties, or reducing alternative routes. We report this gap explicitly and treat it as a gate for the remaining experiments.

\subsection{Experiment 1: Identification and Calibration (RQ1, H1)}

For sampled states and pairs in the tiny-oracle environment we compute the oracle signed effect by clone-state intervention and the estimator score on the same state and pair, and report bias, mean absolute error, Spearman rank correlation, and sign agreement. Three comparisons isolate the contributions: doubly-robust versus plug-in \eqref{eq:dr}; the signed estimand versus the unsigned range statistic, the latter compared against $|{\cdot}|$ of the oracle for fairness; and a sweep over $\varepsilon\in\{0,0.01,0.03,0.05\}$, where $\varepsilon=0$ removes the identification argument entirely and serves as the observational control. We also report the realised forcing rate and its cost in return, so that the price of identification is stated rather than assumed negligible.

\subsection{Experiment 2: Selective Responsiveness (RQ2, H2)}

This experiment measures the paper's central claim directly. Let $W^{(t)}$ collect $\hat\mu_{ij}$ into an influence matrix at evaluation point $t$, and define the normalised movement
\begin{equation}
\Delta^{(t)}=\frac{\big\|W^{(t)}-W^{(t-1)}\big\|_{F}}{\big\|W^{(t-1)}\big\|_{F}+\epsilon},
\qquad
\mathrm{SR}=\frac{\overline{\Delta}_{\mathrm{struct}}}{\overline{\Delta}_{\mathrm{behav}}} .
\label{eq:selectivity}
\end{equation}
A method that separates the tiers should keep $\Delta$ small under behavioural drift, where dependencies are fixed and only policies move, and large after structural shift, giving $\mathrm{SR}\gg1$. Correlation-based weighting should respond to both, giving $\mathrm{SR}\approx1$. We plot both curves on shared axes as the primary figure, and compare against an attention-weighted mean-field baseline. The normalisation in \eqref{eq:selectivity} makes the quantity dimensionless and comparable across methods; an unnormalised variance can be small merely because influence magnitudes are small, which describes a frozen belief rather than a stable one.

\subsection{Experiment 3: Structural Shift Recovery (RQ2)}

We change the influence structure mid-training---reassigning a lane controller, degrading a station, shifting sink demand, or relocating a bottleneck---and report post-shift core F1, recovery latency measured in evaluation intervals with the $H$-step trigger delay accounted for, trigger event times, and reward trajectories. The key comparison is against the ablation without two-timescale separation and residual triggering.

\subsection{Experiment 4: Peripheral Representation and Slot Collapse (RQ3, H3)}

We compare peripheral encoders against a full explicit reference: a single mean-field aggregate, an attention-weighted aggregate, unconstrained soft slots, semantic slots without auxiliary losses, and the full semantic-plus-free routing with \eqref{eq:lb}--\eqref{eq:orth}. Beyond value-prediction error, policy-logit distance from the reference, action agreement, runtime, and memory, we report two collapse diagnostics: slot usage entropy normalised by $\log K$, which detects monopolisation, and the off-diagonal slot cosine similarity of \eqref{eq:orth}, which detects the uniform mode that entropy cannot see. We additionally visualise mean signatures per slot as a heat map; well-separated rows are evidence of specialisation, near-identical rows indicate residual collapse. A recovered-role accuracy against diagnostic ground truth is reported so that the mechanism yields an interpretability result even if returns are unchanged.

\subsection{Experiment 5: Scalability (RQ3)}

Population size is varied over $N\in\{8,16,24,48,96\}$, reporting reward per agent, structural F1, runtime, inference latency, memory, and throughput. Rather than a single operating point we sweep the core budget $k_{\max}$ and plot the reward--compute Pareto frontier, since the practically relevant question is how much return is retained per unit of modelling capacity.

\subsection{Experiment 6: Reciprocity (exploratory diagnostic)}

Everything above measures one direction, $j\to i$. A natural complement asks the reverse: does $j$ react to $i$? We measure it as the information gain from conditioning on the ego when predicting the neighbour's next action,
\begin{equation}
g_{ij}=\underbrace{\mathbb{E}\big[-\log p(a_j^{t+1}\mid s_{ij})\big]}_{\text{ego-blind, shadow state}}-\underbrace{\mathbb{E}\big[-\log p(a_j^{t+1}\mid z_{ij})\big]}_{\text{ego-aware, pair latent}},
\label{eq:reciprocity}
\end{equation}
which is nearly free: the shadow state already excludes ego information and the pair latent already includes it, so only one additional head is required, matched in capacity so that the difference reflects information rather than model size. Computed observationally, $g_{ij}$ is confounded---adjacent agents observe overlapping states---so we evaluate it only on steps where the \emph{ego's} action was itself forced, where $a_i$ is independent of everything by construction and any predictive power is causal.

The two directions define four relation types: strong $j\to i$ with weak $i\to j$ is an indifferent obstacle, to be avoided but not modelled behaviourally; weak $j\to i$ with strong $i\to j$ is a follower; both strong is a strategically coupled pair, which is where the co-adaptation loop responsible for non-stationarity actually resides. We report the resulting scatter and quadrant occupancy. We treat this as diagnostic rather than mechanism: it enters neither core selection nor slot routing here, and we report it partly so that a negative outcome---most pairs falling into a single quadrant, indicating the axis carries no additional information in this environment---is recorded rather than hidden.

\subsection{Baselines and Ablations}

Baselines are pure mean field \cite{yang2018mfmarl}, attention-weighted mean field \cite{hao2023gatmf}, causal mean-field Q-learning \cite{ma2023cmfq} as the closest competing use of interventional weighting, QMIX \cite{rashid2020qmix} as a standard cooperative reference, the unsigned influence measure of \cite{pieroth2024influence} as the closest competing influence estimator, and full explicit local modelling as an upper reference. Ablations remove one mechanism at a time: $\varepsilon$-forcing ($\varepsilon=0$); uncertainty-targeted allocation of the budget \eqref{eq:targeted} (reverting to uniform $\varepsilon$); the doubly-robust correction; the signed estimand (reverting to unsigned magnitude); bias correction \eqref{eq:debias}; controlled forgetting \eqref{eq:inflate}; LCB selection (reverting to the saturating score); the frozen probe \eqref{eq:probe} (reverting to a live-model residual); the matrix detector; semantic slots (unconstrained routing); the auxiliary losses \eqref{eq:lb}--\eqref{eq:orth}; the ego-conditioned terms \eqref{eq:lw}--\eqref{eq:lcon}; and two-timescale separation.

\subsection{Metrics and Statistical Protocol}

Beyond return, structural F1, and throughput, we report the diagnostics that make silent failures visible: ensemble disagreement \eqref{eq:ensemble}, since a degenerate ensemble renders $\hat\sigma$ meaningless; the fraction of updates at which the core budget binds \eqref{eq:budget}, since a saturated budget is not an allocation; slot usage entropy and slot similarity; the pair-specificity ratio; sign agreement against oracle; and the magnitude of the doubly-robust correction term in \eqref{eq:dr}, which tracks outcome-model bias over training.

All results are reported over at least five seeds. Because reported differences in this domain are frequently smaller than seed dispersion, we report bootstrap confidence intervals for each metric and, for pairwise comparisons, a bootstrap distribution over the difference together with the probability that one method exceeds the other; differences whose interval contains zero are described as inconclusive rather than as ordering \cite{agarwal2021precipice}. Evaluation intervals, phase boundaries, warm-up length, and hyperparameters are held constant across compared methods except where they are the object of an ablation.

% ============================================================
\section{Results}

\textcolor{red}{[Author note: the tables and figures below are from the pre-revision run and are retained only as a placeholder. They must be regenerated under the revised method. In particular: (i) core size of $6.0\pm0.0$ reflects a binding cap rather than a learned quantity and should be replaced by the adaptive budget together with its saturation rate; (ii) uncertainty of $0.676\pm0.000$ across variants and tasks indicates a deterministic schedule rather than ensemble disagreement; (iii) temporal variance of order $10^{-8}$ should be replaced by the normalised movement of \eqref{eq:selectivity}; (iv) all reward comparisons require bootstrap intervals, since the reported gaps are smaller than the seed dispersion. Insert the regenerated Experiment 0--5 results here.]}

% ============================================================
\section{Discussion}

Several limitations follow directly from the design and should be weighed against the results.

\emph{Identification is purchased with perturbation.} The interventional estimand is identified because we randomise, and randomisation costs return. The forcing rate is small and annealed, but it is a real intervention in the learning process rather than a free analytical device; the reported cost should be read as part of the method, not as noise.

\emph{Identification covers only part of the data.} With $\varepsilon=0.03$, the randomisation guarantee attaches to roughly three percent of transitions. The outcome model is trained on all of them, so the remaining ninety-seven percent can still imprint spurious associations; the doubly-robust correction pushes back against this, but its influence weight is clipped at $c_{\mathrm{iw}}$ precisely to keep variance finite, and clipping necessarily leaves some bias uncorrected. The honest description is partial identification with extrapolation, not identification outright, and the $\varepsilon$-sweep of Experiment~1 is what makes the residual dependence visible: if calibration does not improve with $\varepsilon$, the mechanism is not doing the work we claim for it.

\emph{Identification is local, not global.} A1 restricts the estimand to short local chains within horizon $H$. Long-range or higher-order dependencies---where $j$ affects $i$ only through a chain of intermediaries, or where two neighbours matter only jointly---are outside the estimator's scope. The framework detects who matters to whom pairwise; it does not perform causal discovery over the interaction graph, and it does not decompose an effect into the paths through which it travels \cite{triantafyllou2024ase}.

\emph{Semantic slots trade flexibility for stability.} Fixing four valence categories buys interpretability and protects against collapse, but presumes that valence is the primary axis of variation. Where the salient distinction is orthogonal to valence---two harmful neighbours differing only in latency---the free slots must carry it, and their capacity is limited. Empirically, partitioning by valence first and specialising within is markedly more effective than global partitioning, because global scaling is dominated by strongly influential relations; but this remains a design choice rather than a derived optimum.

\emph{Computational cost remains the principal weakness.} Explicit pair modelling, an ensemble of outcome models, and slot routing are all more expensive than a single aggregate, and throughput is far below mean-field baselines. The Pareto analysis of Experiment 5 is intended to make this trade-off explicit rather than to conceal it; a practitioner whose bottleneck is wall-clock time should prefer aggregation.

\emph{The convergence statement is narrow.} Proposition 1 concerns the belief recursion given a converging proxy sequence; it does not establish convergence of the coupled proxy--belief--policy system, and the partition rule \eqref{eq:hysteresis} is a discrete map for which we make no stability claim. The error of the estimate is also bounded below by the error of the outcome model, in the sense established for related influence measures \cite{pieroth2024influence}: no post-processing can recover structure that the outcome model failed to represent.

\emph{Controlled forgetting trades one failure mode for another.} Inflation reopens adaptation, but a false detection discards hard-won confidence and transiently shrinks the core. The refractory period and the conjunctive-detector option bound how often this can happen, yet the correct inflation factor remains an empirical choice: too small and rigidity persists, too large and the system forgets a structure that never changed. We report trigger counts and false-alarm behaviour so the setting can be judged rather than assumed.

\emph{The reciprocity axis may carry no signal here.} Experiment~6 is exploratory. In a cooperative resource-flow arena it is plausible that most relations are indifferent obstacles and the strategically coupled quadrant is nearly empty, in which case the second direction adds nothing and should be dropped rather than defended. The causal variant is also sample-hungry: it consumes only those steps where the ego itself was forced, a fraction of an already small budget, which is why we report per-pair sample counts and withhold the estimate rather than reporting a number computed from too few observations.

\emph{Evaluation is single-environment.} All results derive from one arena family plus its tiny-oracle companion. Whether the tier separation transfers to environments with different influence topologies---dense social dilemmas, or settings where structure changes as fast as behaviour, violating A4---is untested.

% ============================================================
\section{Conclusion}

We presented CIG-AMF, which treats structural influence as an interventional estimand and makes it identifiable online by injecting sparse randomised actions, exploiting the fact that in MARL the behaviour policy is known exactly and thus the propensity required by doubly-robust estimation is available in closed form. Retaining sign and horizon structure yields an influence signature rather than a scalar weight, which in turn permits memory slots with fixed functional meaning and a core partition governed by an uncertainty-penalised confidence bound rather than a saturating score. The framework is positioned as a practical bridge between aggregation-based scalability and explicit relational modelling: it does not solve causal discovery in MARL, but it makes a bounded, calibrated, and measurable claim about which neighbours deserve modelling capacity, and it provides diagnostics under which failures of that claim are visible rather than silent.

\section*{Acknowledgment}

% ============================================================
\begin{thebibliography}{00}

\bibitem{yang2018mfmarl} Y. Yang, R. Luo, M. Li, M. Zhou, W. Zhang, and J. Wang, ``Mean field multi-agent reinforcement learning,'' in \textit{Proc. 35th Int. Conf. Machine Learning}, 2018.

\bibitem{busoniu2008survey} L. Busoniu, R. Babuska, and B. De Schutter, ``A comprehensive survey of multiagent reinforcement learning,'' \textit{IEEE Trans. Syst., Man, Cybern. C}, vol. 38, no. 2, pp. 156--172, 2008.

\bibitem{papoudakis2019nonstationarity} G. Papoudakis, F. Christianos, A. Rahman, and S. V. Albrecht, ``Dealing with non-stationarity in multi-agent deep reinforcement learning,'' arXiv:1906.04737, 2019.

\bibitem{hao2023gatmf} Q. Hao, W. Huang, T. Feng, J. Yuan, and Y. Li, ``GAT-MF: Graph attention mean field for very large scale multi-agent reinforcement learning,'' in \textit{Proc. ACM SIGKDD Conf. Knowledge Discovery and Data Mining}, 2023.

\bibitem{wang2024amf} X. Wang, Y. Wang, Y. Wang, and X. Zhao, ``Adaptive mean field multi-agent reinforcement learning,'' \textit{Information Sciences}, vol. 669, 2024.

\bibitem{ma2023cmfq} H. Ma, C. Peng, W. Li, T. Li, L. Gui, L. Gao, and S. Gyawali, ``Causal mean field multi-agent reinforcement learning,'' in \textit{Int. Joint Conf. Neural Networks}, 2023.

\bibitem{jaques2019influence} N. Jaques, A. Lazaridou, E. Hughes, C. Gulcehre, P. A. Ortega, D. Strouse, J. Z. Leibo, and N. de Freitas, ``Social influence as intrinsic motivation for multi-agent deep reinforcement learning,'' in \textit{Proc. 36th Int. Conf. Machine Learning}, 2019.

\bibitem{du2024scic} X. Du, Y. Ye, P. Zhang, Y. Yang, M. Chen, and T. Wang, ``Situation-dependent causal influence-based cooperative multi-agent reinforcement learning,'' in \textit{Proc. AAAI Conf. Artificial Intelligence}, 2024.

\bibitem{yu2026magic} H. Yu, L. Wang, J. Cong, S. Wang, and C. Liu, ``MAGIC: Multi-step advantage-gated causal influence for multi-agent reinforcement learning,'' arXiv:2605.01805, 2026.

\bibitem{pieroth2024influence} F. R. Pieroth, K. Fitch, and L. Belzner, ``Detecting influence structures in multi-agent reinforcement learning,'' in \textit{Proc. 41st Int. Conf. Machine Learning}, PMLR 235, 2024.

\bibitem{triantafyllou2024ase} S. Triantafyllou, A. Sukovic, D. Mandal, and G. Radanovic, ``Agent-specific effects: A causal effect propagation analysis in multi-agent MDPs,'' in \textit{Proc. 41st Int. Conf. Machine Learning}, 2024.

\bibitem{foerster2018coma} J. N. Foerster, G. Farquhar, T. Afouras, N. Nardelli, and S. Whiteson, ``Counterfactual multi-agent policy gradients,'' in \textit{Proc. AAAI Conf. Artificial Intelligence}, 2018.

\bibitem{wolpert2002difference} D. H. Wolpert and K. Tumer, ``Optimal payoff functions for members of collectives,'' \textit{Advances in Complex Systems}, vol. 4, no. 2/3, pp. 265--279, 2002.

\bibitem{tumer2007aristocrat} K. Tumer and A. Agogino, ``Distributed agent-based air traffic flow management,'' in \textit{Proc. 6th Int. Conf. Autonomous Agents and Multiagent Systems}, 2007.

\bibitem{rashid2020qmix} T. Rashid, M. Samvelyan, C. Schroeder de Witt, G. Farquhar, J. Foerster, and S. Whiteson, ``Monotonic value function factorisation for deep multi-agent reinforcement learning,'' \textit{J. Machine Learning Research}, vol. 21, no. 178, pp. 1--51, 2020.

\bibitem{boehmer2020dcg} W. B\"{o}hmer, V. Kurin, and S. Whiteson, ``Deep coordination graphs,'' in \textit{Proc. 37th Int. Conf. Machine Learning}, 2020.

\bibitem{wang2022casec} T. Wang, L. Zeng, W. Dong, Q. Yang, Y. Yu, and C. Zhang, ``Context-aware sparse deep coordination graphs,'' in \textit{Int. Conf. Learning Representations}, 2022.

\bibitem{wang2020roma} T. Wang, H. Dong, V. Lesser, and C. Zhang, ``ROMA: Multi-agent reinforcement learning with emergent roles,'' in \textit{Proc. 37th Int. Conf. Machine Learning}, 2020.

\bibitem{wang2021rode} T. Wang, T. Gupta, A. Mahajan, B. Peng, S. Whiteson, and C. Zhang, ``RODE: Learning roles to decompose multi-agent tasks,'' in \textit{Int. Conf. Learning Representations}, 2021.

\bibitem{zeng2023sird} X. Zeng, H. Peng, and A. Li, ``Effective and stable role-based multi-agent collaboration by structural information principles,'' in \textit{Proc. AAAI Conf. Artificial Intelligence}, 2023.

\bibitem{hu2024acorm} Z. Hu, Z. Zhang, H. Li, C. Chen, H. Ding, and Z. Wang, ``Attention-guided contrastive role representations for multi-agent reinforcement learning,'' in \textit{Int. Conf. Learning Representations}, 2024.

\bibitem{seitzer2021cai} M. Seitzer, B. Sch\"{o}lkopf, and G. Martius, ``Causal influence detection for improving efficiency in reinforcement learning,'' in \textit{Advances in Neural Information Processing Systems}, 2021.

\bibitem{jiang2016dr} N. Jiang and L. Li, ``Doubly robust off-policy value evaluation for reinforcement learning,'' in \textit{Proc. 33rd Int. Conf. Machine Learning}, 2016.

\bibitem{thomas2016hcope} P. Thomas and E. Brunskill, ``Data-efficient off-policy policy evaluation for reinforcement learning,'' in \textit{Proc. 33rd Int. Conf. Machine Learning}, 2016.

\bibitem{lakshminarayanan2017ensembles} B. Lakshminarayanan, A. Pritzel, and C. Blundell, ``Simple and scalable predictive uncertainty estimation using deep ensembles,'' in \textit{Advances in Neural Information Processing Systems}, 2017.

\bibitem{fedus2022switch} W. Fedus, B. Zoph, and N. Shazeer, ``Switch Transformers: Scaling to trillion parameter models with simple and efficient sparsity,'' \textit{J. Machine Learning Research}, vol. 23, no. 120, pp. 1--40, 2022.

\bibitem{locatello2020slotattention} F. Locatello, D. Weissenborn, T. Unterthiner, A. Mahendran, G. Heigold, J. Uszkoreit, A. Dosovitskiy, and T. Kipf, ``Object-centric learning with slot attention,'' in \textit{Advances in Neural Information Processing Systems}, 2020.

\bibitem{borkar2008sa} V. S. Borkar, \textit{Stochastic Approximation: A Dynamical Systems Viewpoint}. Hindustan Book Agency, 2008.

\bibitem{kingma2015adam} D. P. Kingma and J. Ba, ``Adam: A method for stochastic optimization,'' in \textit{Int. Conf. Learning Representations}, 2015.

\bibitem{page1954cusum} E. S. Page, ``Continuous inspection schemes,'' \textit{Biometrika}, vol. 41, no. 1/2, pp. 100--115, 1954.

\bibitem{zhang2024corep} J. Zhang, J. Wang, H. Hu, Y. Chen, C. Fan, and C. Zhang, ``Tackling non-stationarity in reinforcement learning via causal-origin representation,'' in \textit{Proc. 41st Int. Conf. Machine Learning}, 2024.

\bibitem{nguyen2011baipomdp} T. H. Nguyen, D. Hsu, W. S. Lee, T. L. Leong, L. P. Kaelbling, T. Lozano-Perez, and A. H. Grant, ``Bayes-adaptive interactive POMDPs,'' in \textit{Proc. AAAI Conf. Artificial Intelligence}, 2011.

\bibitem{zintgraf2020varibad} L. Zintgraf, K. Shiarlis, M. Igl, S. Schulze, Y. Gal, K. Hofmann, and S. Whiteson, ``VariBAD: A very good method for Bayes-adaptive deep RL via meta-learning,'' in \textit{Int. Conf. Learning Representations}, 2020.

\bibitem{agarwal2021precipice} R. Agarwal, M. Schwarzer, P. S. Castro, A. Courville, and M. G. Bellemare, ``Deep reinforcement learning at the edge of the statistical precipice,'' in \textit{Advances in Neural Information Processing Systems}, 2021.

\bibitem{leibo2021meltingpot} J. Z. Leibo, E. Due\~{n}ez-Guzm\'{a}n, A. Vezhnevets, J. P. Agapiou, P. Sunehag, R. Koster, J. Matyas, C. Beattie, I. Mordatch, and T. Graepel, ``Scalable evaluation of multi-agent reinforcement learning with Melting Pot,'' in \textit{Proc. 38th Int. Conf. Machine Learning}, 2021.

\bibitem{ellis2023smacv2} B. Ellis, J. Cook, S. Moalla, M. Samvelyan, M. Sun, A. Mahajan, J. N. Foerster, and S. Whiteson, ``SMACv2: An improved benchmark for cooperative multi-agent reinforcement learning,'' in \textit{Advances in Neural Information Processing Systems}, 2023.

\bibitem{suarez2019neuralmmo} J. Suarez, Y. Du, P. Isola, and I. Mordatch, ``Neural MMO: A massively multiagent game environment for training and evaluating intelligent agents,'' arXiv:1903.00784, 2019.

\bibitem{carroll2019overcooked} M. Carroll, R. Shah, M. K. Ho, T. L. Griffiths, S. A. Seshia, P. Abbeel, and A. Dragan, ``On the utility of learning about humans for human-AI coordination,'' in \textit{Advances in Neural Information Processing Systems}, 2019.

\bibitem{zaheer2017deepsets} M. Zaheer, S. Kottur, S. Ravanbakhsh, B. Poczos, R. Salakhutdinov, and A. Smola, ``Deep Sets,'' in \textit{Advances in Neural Information Processing Systems}, 2017.

\bibitem{lee2019settransformer} J. Lee, Y. Lee, J. Kim, A. R. Kosiorek, S. Choi, and Y. W. Teh, ``Set Transformer: A framework for attention-based permutation-invariant neural networks,'' in \textit{Proc. 36th Int. Conf. Machine Learning}, 2019.

\end{thebibliography}

\end{document}
